%% start of file `template-zh.tex'.
%% Copyright 2006-2013 Xavier Danaux (xdanaux@gmail.com).
%
% This work may be distributed and/or modified under the
% conditions of the LaTeX Project Public License version 1.3c,
% available at http://www.latex-project.org/lppl/.


\documentclass[12pt,a4paper,sans]{moderncv}   % possible options include font size ('10pt', '11pt' and '12pt'), paper size ('a4paper', 'letterpaper', 'a5paper', 'legalpaper', 'executivepaper' and 'landscape') and font family ('sans' and 'roman')

% moderncv 主题
% \moderncvstyle{classic}                        % 选项参数是 ‘casual’, ‘classic’, ‘oldstyle’ 和 ’banking’
% \moderncvcolor{blue}                          % 选项参数是 ‘blue’ (默认)、‘orange’、‘green’、‘red’、‘purple’ 和 ‘grey’
\moderncvstyle{banking}
% \moderncvtheme[blue]{classic}
% \nopagenumbers{}                             % 消除注释以取消自动页码生成功能

\pagestyle{fancy}
\fancyhf{}
\fancyfoot[L]{\the\year/\the\month}
\fancyfoot[C]{\small \textit{Haoyang Lu}}
\fancyfoot[R]{\thepage}

% 字符编码
\usepackage[utf8]{inputenc}                   % 替换你正在使用的编码
% \usepackage{CJKutf8}
% \usepackage{xltxtra}
\usepackage{fontspec}
\usepackage{xunicode}
\usepackage{xeCJK}
\setmainfont{Merriweather}
\setsansfont{Source Sans 3}
% \setsansfont{Helvetica Neue-Light}
\setmonofont{Courier New}
\setCJKmainfont{Noto Serif CJK SC}
\setCJKsansfont{Noto Sans CJK SC}
\setCJKmonofont{Noto Sans Mono CJK SC}

% 调整页边距
% \usepackage[scale=0.9]{geometry}
\usepackage[margin=0.5in]{geometry}
\recomputelengths
\setlength{\hintscolumnwidth}{2cm}           % 如果你希望改变日期栏的宽度

% 对列表各项逆序编号（用于对文章进行编号）
\usepackage[itemsep=2pt]{etaremune}

% 个人信息
\name{Haoyang}{Lu}
% \title{简历题目 (可选项)}                     % 可选项、如不需要可删除本行
\address{Russell Square House, 10-12 Russell Square}{London,  WC1B 5EH}{UK}
% \phone[mobile]{+44-7432843788}              % 可选项、如不需要可删除本行
% \phone[fixed]{+2~(345)~678~901}               % 可选项、如不需要可删除本行
% \phone[fax]{+3~(456)~789~012}                 % 可选项、如不需要可删除本行
\email{haoyang.lu@ucl.ac.uk}                    % 可选项、如不需要可删除本行
\social[github]{hy-lu}
\social[researchgate]{Haoyang-Lu}
% \social[linkedin]{haoyang-lu}
% \homepage{www.xialongli.com}                  % 可选项、如不需要可删除本行
% \extrainfo{附加信息 (可选项)}                 % 可选项、如不需要可删除本行
% \photo[64pt][0.4pt]{picture}                  % ‘64pt’是图片必须压缩至的高度、‘0.4pt‘是图片边框的宽度 (如不需要可调节至0pt)、’picture‘ 是图片文件的名字;可选项、如不需要可删除本行
% \quote{引言(可选项)}                          % 可选项、如不需要可删除本行

% 显示索引号;仅用于在简历中使用了引言
%\makeatletter
%\renewcommand*{\bibliographyitemlabel}{\@biblabel{\arabic{enumiv}}}
%\makeatother

% 分类索引
%\usepackage{multibib}
%\newcites{book,misc}{{Books},{Others}}
%----------------------------------------------------------------------------------
%            内容
%----------------------------------------------------------------------------------
\begin{document}
% \begin{CJK}{UTF8}{gbsn}                       % 详情参阅CJK文件包
\maketitle

\section{Employment}
\cventry{2024 -- Present}{Research Fellow in Computational Psychiatry}{Max Planck UCL Centre for Computational Psychiatry and Ageing Research}{London, UK}{PI: Dr. Quentin Huys}{Develop and apply reliable behavioral tasks and computational models for randomized controlled trials to investigate the reinforcement learning mechanisms of two types of antidepressants in people with depression (RELMED).}

\cventry{2022 -- 2024}{Postdoctoral Research Fellow}{School of Psychological and Cognitive Sciences, Peking University}{Beijing, China}{PI: Dr. Hang Zhang}{Conduct large-scale online experiments and use computational models to investigate the cognitive processes and influencing factors involved in the formation and updating of superstitious beliefs and the relationship with clinical/personality traits.}

\section{Education}
% \cventry{years}{degree/job title}{institution/employer} {localization}{optional: grade/...}{optional: comment/job description}
\cventry{2016 -- 2022}{PhD in Integrated Life Sciences (Psychology)}{Peking University}{Beijing, China}{}{
  \begin{itemize}
    \item Dissertation Title: \emph{Inference and Decision-making in People with Autism Spectrum Disorder or Broader Autism Phenotype}
    \item Advisors: Dr. Li Yi, Dr. Hang Zhang
  \end{itemize}
}  % 第3到第6编码可留白
\cventry{2012 -- 2016}{BSc in Psychology}{Sun Yat-Sen University}{Guangzhou, China}{}{
}

% \section{Thesis}
% % %to typeset lines with a hint on the left: 
% % \cvline{leftmark}{text} or \cvitem{leftmark}{text}
% \cvitem{Title}{\emph{Inference and Decision-making in People with Autism Spectrum Disorder or Broader Autism Phenotype}}
% \cvitem{Advisors}{Li Yi, Hang Zhang}
% \cvitem{Introduction}{\small The core features of autism spectrum disorders (ASD) are believed to be pertinent to how individuals interact with and sample the world. Therefore, it is crucial to understand the outcome of atypical information sampling in ASD. In this interdisciplinary project, I worked with two advisors, Prof.Li Yi, who specializes in child psychopathology and Prof.Hang Zhang in computational cognitive science. I conducted a series of studies on information sampling in both autistic children and adults with a broader autism phenotype. Through behavioral tasks, eye-tracking/pupillometry, and computational modeling, I uncovered distinct behavioral, attentional, and cognitive processes in autistic individuals compared to neurotypical people for both instrumental and non-instrumental information sampling.}

\section{Publications}
% \subsection{Preprints or under review}
% \newcommand{\Revision}{\textit{Under revision}}
\newcommand{\Review}{\textit{Under review}}
\newcommand{\Submitted}{\textit{Submitted}}
\newcommand{\Manuscript}{\textit{Preparing manuscript}}
\newcommand{\Preprint}{\textit{Preprint}}
\newcommand{\CS}{*} % corresponding author
\newcommand{\CF}{\textsuperscript{\#}} % co-first author
\newcommand{\ME}{\textbf{Lu, H.}}

% \CS corresponding author, \CF co-first author.

% \begin{enumerate}

% \end{enumerate}

\subsection{Peer-reviewed journal articles}
\input{journal}
\subsection{Conference abstracts}
\input{conference}

\section{Professional membership and service}
\cvitem{Membership}{Society for Research in Child Development; International Society for Autism Research; Cognitive Science Society; Society for Neuroeconomics}
\cvitem{Ad hoc reviewer}{Autism Research; Journal of Autism and Developmental Disorders; eLife; OpenMind; Computational Psychiatry Conference}

\section{Professional skills and languages}
% \cvitemwithcomment{}{fNIRS, fMRI, EEG}{Beginner}
\cvitem{Research}{Eye-tracking, Hierarchical Bayesian modeling, \LaTeX}
\cvitem{Statistics}{Linear mixed model, Bayesian statistics, Generalized additive model, Machine learning methods}
\cvitem{Programming}{R, Julia, JavaScript, MATLAB, Stan, Python, Git}
\cvitem{Languages}{Chinese, English}

\clearpage
% \end{CJK}
\end{document}


%% 文件结尾 `template-zh.tex'.
